\let\negmedspace\undefined
\let\negthickspace\undefined
\documentclass[journal]{IEEEtran}
\usepackage[a5paper, margin=10mm, onecolumn]{geometry}
\usepackage{lmodern} % Ensure lmodern is loaded for pdflatex
\usepackage{tfrupee} % Include tfrupee package

\setlength{\headheight}{1cm} % Set the height of the header box
\setlength{\headsep}{0mm}     % Set the distance between the header box and the top of the text

\usepackage{gvv-book}
\usepackage{gvv}
\usepackage{cite}
\usepackage{amsmath,amssymb,amsfonts,amsthm}
\usepackage{algorithmic}
\usepackage{graphicx}
\graphicspath{{./figs/}}
\usepackage{textcomp}
\usepackage{xcolor}
\usepackage{txfonts}
\usepackage{listings}
\usepackage{enumitem}
\usepackage{mathtools}
\usepackage{gensymb}
\usepackage{comment}
\usepackage[breaklinks=true]{hyperref}
\usepackage{tkz-euclide} 
\usepackage{listings}
\usepackage{gvv}                                        
\def\inputGnumericTable{}                                 
\usepackage[latin1]{inputenc}                                
\usepackage{color}                                            
\usepackage{array}                                            
\usepackage{longtable}                                       
\usepackage{calc}                                             
\usepackage{multirow}                                         
\usepackage{hhline}                                           
\usepackage{ifthen}                                           
\usepackage{lscape}
\usepackage{circuitikz}
\tikzstyle{block} = [rectangle, draw, fill=blue!20, 
text width=4em, text centered, rounded corners, minimum height=3em]
\tikzstyle{sum} = [draw, fill=blue!10, circle, minimum size=1cm, node distance=1.5cm]
\tikzstyle{input} = [coordinate]
\tikzstyle{output} = [coordinate]


\begin{document}
	
	\bibliographystyle{IEEEtran}
	\vspace{3cm}
	
	\title{1.5.4}
	\author{EE25BTECH11016 - Taraka Abhinav}
	\maketitle
	{\let\newpage\relax\maketitle}
	
	\renewcommand{\thefigure}{\theenumi}
	\renewcommand{\thetable}{\theenumi}
	\setlength{\intextsep}{10pt} % Space between text and floats
	
	
	\numberwithin{equation}{enumi}
	\numberwithin{figure}{enumi}
	\renewcommand{\thetable}{\theenumi}

\textbf{Question:}\\
A circle has its center at $(4,4)$. If one end of a diameter is $(4,0)$, then find the coordinates of the other end.\\
\solution\\
According to the question,  
Consider the coordinates as vectors:

\begin{align*}
A &= \myvec{x\\y}, \quad B = \myvec{4\\0}, \quad C = \myvec{4\\4}
\end{align*}

The midpoint of the diameter is the center. Hence,
\begin{align*}
\frac{A+B}{2} &= C
\end{align*}

Multiplying both sides by $2$,
\begin{align*}
A + B &= 2C
\end{align*}

Substituting,
\begin{align*}
\myvec{x\\y} + \myvec{4\\0} &= 2\myvec{4\\4}
\end{align*}

Simplifying,
\begin{align*}
\myvec{x\\y} &= 2\myvec{4\\4} - \myvec{4\\0} \\
&= \myvec{8\\8} - \myvec{4\\0} \\
&= \myvec{4\\8}
\end{align*}

Now writing as a linear system,
\begin{align*}
\myvec{1 & 0\\0 & 1}\myvec{x\\y} &= \myvec{4\\8}
\end{align*}

The augmented matrix is,
\begin{align*}
\myvec{1 & 0 & 4\\0 & 1 & 8}
\end{align*}

Row operations give the same matrix :
\begin{align*}
\myvec{1 & 0 & 4\\0 & 1 & 8}
\end{align*}

Thus,
\begin{align*}
\operatorname{rank}\myvec{1 & 0\\0 & 1} 
&= \operatorname{rank}\myvec{1 & 0 & 4\\0 & 1 & 8} = 2
\end{align*}

So the unique solution is
\begin{align*}
\myvec{x\\y} &= \myvec{4\\8}
\end{align*}

\begin{align*}
\therefore \text{The other end of the diameter is }(4,8).
\end{align*}

\begin{figure}[H]
    \centering
    \includegraphics[width=0.5\linewidth]{figs/Figure_1.png}
    \label{figure_1}
\end{figure}

\end{document}